% !TEX root = ../main.tex
% !TEX program = xelatex
\newpage
\section{数据预处理与探索性分析}

\subsection{数据预处理}

本文数据来自赛题附件，涵盖负载、光伏、电价及不同时间发布的光伏预报。各附件的时间写法和排列方式有所不同，功率数据也需要转换为调度所用的电量。为使这些数据能够准确对应，我们进行了以下处理。

\noindent\textbf{1. 检查缺失与异常记录。}检查空单元格、非数值项和极端取值。原表已有记录未发现空单元格，但按实际时间排列后，全年实测数据缺少1月1日00:00—00:10这一段，处理方法见第6.5节。夜间光伏为零属于正常情况，予以保留；较大或较小的观测值经原表核对后，也不因数值极端而直接删除。

\noindent\textbf{2. 统一时间与区间对应关系。}将各行数据用于其标注时刻开始的10分钟区间。例如，首行“00:10”对应00:10—00:20，末行“00:00+1”对应次日00:00—00:10。问题一中，每天各时段的负载、光伏出力和电价均相同，因此将末行零点数据放到当天开头；全年连续数据则按实际日期衔接，不将次日零点移回当天。描述性统计仍按原表日期分组。

\noindent\textbf{3. 整理数据并区分预报批次。}统一日期、时间格式，将各附件整理为便于按时刻匹配的记录。光伏预报同时保留“何时发布”和“预测哪个时刻”两项信息，避免新预报覆盖旧预报，也便于判断决策时有哪些预报可用。超出实测覆盖范围的预报予以保留，但不参与误差比较。

\noindent\textbf{4. 转换电量并核对数据。}每个时段为10分钟，即六分之一小时，因此将负载和光伏功率乘以六分之一，得到对应时段的电量，即 \(E=P\Delta t\)，其中 \(\Delta t=1/6\,\mathrm h\)。最后核对整理前后的数值与时间对应关系。全年统计用于分析数据特征，逐日决策则仅使用当时已经获得的信息。

\newpage
\subsection{探索性分析}

经过整理与分析，得到以下四点发现。

\subsubsection*{（1）光伏富余与用电缺口存在时间错位，储能需要协调全天安排}

\begin{figure}[H]
    \centering
    \includegraphics[width=0.8\textwidth]{fig/tsxfx_1.png}
    \caption{代表日供需与电价的时间错配}
    \label{fig:tsxfx_1}
\end{figure}

图\ref{fig:tsxfx_1}显示，代表日午间光伏能够覆盖负载并产生富余，傍晚光伏回落后，外网需要补足的电量明显增加，且这一时段与较高电价有所重合。全年数据中，有330天出现过光伏功率高于负载的情况，说明光伏富余并非个别现象。

因此，\textbf{全天总电量是否充足，不能代替各时段的供需判断}。储能可以将富余光伏或低价购电留到后续使用，但实际能够转移多少，还受到容量、充放电功率和效率限制。这也是需要统筹全天购电与充放电，而不能逐时独立决策的原因。

\subsubsection*{（2）负载与光伏的年内变化不同步，冬季补电压力未必更小}

\begin{figure}[H]
    \centering
    \includegraphics[width=0.8\textwidth]{fig/tsxfx_2.png}
    \caption{负载、光伏与电价的年内变化}
    \label{fig:tsxfx_2}
\end{figure}

图\ref{fig:tsxfx_2}表明，本年度负载和光伏均有明显的年内变化，但二者并不同步。夏季负载较高，光伏也能提供较多电量；冬季负载虽有所回落，光伏下降却更加明显。以月平均功率差衡量，12月负载减去光伏约为 \(3177\,\mathrm{kW}\)，接近5月的两倍。

这说明，\textbf{外网补电压力取决于负载与光伏的共同变化，不能只看用电量高低}。同时，电价也存在年内起伏。不同月份的供需缺口和购电条件均有差异，因此需要随近期数据更新判断，不能将代表日的安排直接复制到全年。

\subsubsection*{（3）负载具有星期规律，光伏与电价的日际变化更强}

\begin{figure}[H]
    \centering
    \includegraphics[width=0.8\textwidth]{fig/tsxfx_3.png}
    \caption{负载的周重复特征与电价的跨日波动}
    \label{fig:tsxfx_3}
\end{figure}

按实际星期分组后，周五、周六的平均负载比其余五天低约36.8\%；负载与一周前数据的相关性也高于与前一天数据的相关性。因此，\textbf{同一星期位置的历史负载更有参考价值}。

光伏的情况不同。原始曲线在不同日期之间较为相似，但扣除各月相同时刻的平均日内轮廓后，其一周相关性接近于零。这表明，曲线相似主要来自白天发电、夜间为零的共同形状，不能据此认定光伏具有稳定的星期周期。因此，负载和光伏应分别预测，避免将负载的星期规律直接用于光伏。

图\ref{fig:tsxfx_3}中的电价百分位范围还显示，即使处于相同的日内时刻，不同日期的电价仍可能有较大差异。历史价格规律可以提供参考，却不能代替对未来价格的判断。

\subsubsection*{（4）日内新预报总体更准确，但精度提升不等于同比例节费}

\begin{figure}[H]
    \centering
    \includegraphics[width=0.8\textwidth]{fig/tsxfx_4.png}
    \caption{光伏预报提前量与日内更新的误差比较}
    \label{fig:tsxfx_4}
\end{figure}

将不同发布时间的光伏预报与相同目标时刻的实测值匹配后，图\ref{fig:tsxfx_4}显示，预报提前时间越长，平均绝对误差总体越大。在实际光伏大于零的4177组配对样本中，使用目标时刻前最新日内预报后，平均绝对误差由约 \(378.8\,\mathrm{kW}\) 降至 \(116.3\,\mathrm{kW}\)，下降约69\%。但仍有约17.4\%的样本在更新后误差增大，说明\textbf{新预报总体更准，却并非每次都更准}。

预报改善为调整购电提供了依据，但不意味着费用会按相同比例下降。部分供需偏差可以由储能缓冲，修改订单又会产生费用。因此，是否利用新预报改单，应比较调整前后的总费用，综合考虑普通购电、紧急购电和调整成本。这也意味着，问题三需要同时评价预报精度与实际节费效果。
